% SHARD-ID: NB-01-PROGRAMME-MAP
% SHARD-TITLE: Programme Map
% SHARD-SUMMARY: Lists the research axes the notebook covers and the order in which their foundations must be laid.
% SHARD-SUMMARY: Asserts no connections between axes that have not been established.
% SHARD-KEYWORDS: programme, research axes, scope, roadmap

\section{Programme Map}
\label{sec:programme-map}

\subsection*{Purpose of this shard}

This shard says what the notebook is \emph{about}, so that a reader --- human or agent --- can tell
within a minute whether a new idea belongs here and what it would have to connect to. It is a map
of intentions, not of results. Anything claimed as a result belongs in a later shard with its gates
named.

\subsection*{Axes}

\status{open} No axis has been opened yet. Each axis entry, once added, states the question, the
sources that ground it, and the first thing that must be settled before work on it can proceed.

\subsection*{Rules for this map}

\begin{enumerate}
  \item An axis is added when there is a concrete question, not when there is an interest.
  \item Connections between axes are stated only when a derivation or a run supports them. A
        suspected connection is written as \evI{} or \evQ{} and lives in \texttt{ideas/} until it
        is earned.
  \item Each axis names the conventions it depends on. Conflicting conventions between axes are
        resolved in \texttt{CONVENTIONS.md} before shared code is written.
  \item An axis that is dropped is not deleted: it is marked dropped, with the reason, and the
        attempt is recorded in \texttt{attic/}.
\end{enumerate}

\subsection*{What must exist before technical shards}

\begin{itemize}
  \item Registered primary sources for the axis, with local PDFs and hashes.
  \item Unit, sign and normalisation conventions for the quantities involved.
  \item Glossary entries for the terms the axis introduces.
  \item At least one validation target: a known analytic limit, a published number, or a figure to
        reproduce, against which our first computation can be checked.
\end{itemize}
